\chapter{Your first part: a bracket in twenty minutes}

This chapter is one complete, reproducible exercise. Work through it and you will use the most important
modelling chain in the app: \textbf{Sketch → Pad → Fillet → Hole → Save}. When you have finished, you can
already build simple parts with MobileCAD.

\begin{practice}[Goal]
A rectangular bracket 60 × 40 mm and 12 mm thick, with R6 fillets at the four corners, a ⌀10 through hole in
the centre and one ⌀5 mounting hole on each side.
\end{practice}

\section{Step 1: Create a document}

On launch you get an empty document called “Unnamed”. If a document is already open, create a new one with
\menupath{File → New} (\cmd{Std\_New}). Save it once and give it a name:

\begin{enumerate}
  \item \menupath{File → Save} (\cmd{Std\_Save});
  \item in the system file picker choose a folder (Documents, for example) and enter \file{bracket} as the
        file name;
  \item tap OK. The document name in the top bar becomes \file{bracket}.
\end{enumerate}

\begin{tip}[Why save first?]
MobileCAD autosaves at intervals and restores from a recovery copy after a crash, but saving first gives this
session a definite home and makes it easy to use Save a copy later to compare design variants.
\end{tip}

\section{Step 2: Create a sketch and choose a datum plane}

\begin{enumerate}
  \item Make sure the \ui{Part} tab is selected at the bottom, then tap \ui{Sketch} in the top toolstrip;
  \item in the tool panel tap \ui{New Sketch} (command \cmd{Sketcher\_NewSketch});
  \item a \ui{Sketch plane} chooser appears: \ui{Top (XY)}, \ui{Front (XZ)}, \ui{Right (YZ)}. Choose
        \ui{Top (XY)}.
\end{enumerate}

\screenshot{ui-sketch-plane-zh.jpg}{Choosing a datum plane when you create a sketch. Later you can use the “Attachment” property to move the sketch onto a face of the body or another datum plane.}

\section{Step 3: Draw a rectangle}

In sketch mode the toolstrip changes to sketch tools: four toolsets, \ui{Draw}, \ui{Constrain}, \ui{Modify}
and \ui{View}, plus \ui{Construction: off}, \ui{Constraints} (the constraints list), the common
\ui{Polyline}, \ui{Rectangle}, \ui{Circle}, \ui{3-point arc}, \ui{Dimension}, \ui{Trim}, \ui{Select} and
\ui{Delete}, and at the far right ✓ (Leave sketch) and ✗ (Cancel sketch).

\begin{enumerate}
  \item Tap \ui{Rectangle};
  \item tap once near the origin (the first corner), move your finger and tap again (the opposite corner). Do
        not aim for exact sizes; the shape only has to be roughly right — the next step locks it down with
        dimensions.
\end{enumerate}

\screenshot{ui-sketch-tools-zh.jpg}{The sketch tool panel: the “Sketch” toolset lists every sketch command in sections for geometry, constraints, editing and view; once you enter a sketch, the top toolstrip switches to the four sketch groups Draw / Constrain / Modify / View.}

\section{Step 4: Add dimensions and make the sketch fully constrained}

\begin{enumerate}
  \item Tap \ui{Dimension} (\cmd{Sketcher\_Dimension}), tap the rectangle's two horizontal edges or two corners
        in turn, drag to a good position and tap to place the dimension;
  \item the Numpad appears once the dimension is placed; type \file{60} and tap ✓ to confirm;
  \item dimension the other direction in the same way and enter \file{40};
  \item watch the status in the bottom-left corner: when it reads \ui{Fully constrained}, every degree of
        freedom in the sketch is constrained. If it reads “N degrees of freedom”, constraints are still
        missing and the model can still be dragged out of shape.
\end{enumerate}

\begin{tip}[Why does full constraint matter?]
Fully constrained means the geometry is mathematically unique. Later features (Pad, Pocket) depend on this
sketch, and an under-constrained sketch can “jump” when a parameter changes, which makes features fail. Make it
a habit: \textbf{do not leave a sketch until the status bar reads “Fully constrained”}.
\end{tip}

\section{Step 5: Leave the sketch}

Tap ✓ at the far right of the toolstrip (its hint is “Leave sketch”). The view returns to 3D and the feature
tree gains a \ui{Sketch}, nested inside a new \ui{Body} — that is the Part Design rule: every feature belongs
to a body.

\section{Step 6: Pad}

\begin{enumerate}
  \item Select the sketch you just drew in the tree;
  \item tap \ui{Features} → \ui{Pad} in the toolstrip (\cmd{PartDesign\_Pad});
  \item the feature dialog opens with \ui{Profile} already filled in. Tap the \ui{Length} row and type
        \file{12} on the Numpad;
  \item want to see the extrusion direction straight away? Flip the \ui{Reversed} switch;
  \item tap ✓ in the title bar to confirm. The block appears.
\end{enumerate}

\begin{warn}[What if the feature fails?]
If a red error appears below the dialog (for example “cannot find the geometry on the base feature”), the
profile has a problem: it is not closed, it self-intersects or it has extra lines. Cancel, go back to the
sketch, check for conflicts in the \ui{Constraints} list, remove the extra segments with \ui{Trim} and try
again.
\end{warn}

\section{Step 7: Fillet the corners}

\begin{enumerate}
  \item Tap \ui{Features} → \ui{Fillet} in the toolstrip (\cmd{PartDesign\_Fillet});
  \item the \ui{Base} row already points at the current body. Tap \ui{Edges} (or tap the four vertical edges
        directly on the model);
  \item enter \file{6} in \ui{Radius};
  \item tap ✓. The four vertical edges become R6 fillets.
\end{enumerate}

\begin{tip}[The edge table in fillets and chamfers]
The Fillet and Chamfer dialogs let you give each edge its own radius: every edge you add gets a row, and you
can tap the value in a row to change it. To change them all at once, use the \ui{Radius for all} button.
\end{tip}

\section{Step 8: The central through hole}

There are two ways to do this; the first is recommended:

\subsection{Method A: the Hole feature}

\begin{enumerate}
  \item Select the \textbf{top face} first (tap the upper surface);
  \item tap \ui{Features} → \ui{Hole} (\cmd{PartDesign\_Hole});
  \item in the dialog set \ui{Diameter} to \file{10} and \ui{Depth type} to \ui{Through all};
  \item tap ✓.
\end{enumerate}

\begin{tip}[Standard threads and counterbores]
The \ui{Threaded} switch on the Hole feature generates threaded holes to ISO standards (choose thread type,
size and class). \ui{Hole type} offers counterbore, countersink, counterdrill and other shapes, and shows the
matching diameter, depth and angle parameters automatically.
\end{tip}

\subsection{Method B: a sketch plus a Pocket}

Create a new sketch on the top face → draw a ⌀10 circle → leave the sketch → select the sketch and tap
\ui{Pocket} (\cmd{PartDesign\_Pocket}) → set the depth to “Through all”. Both methods give the same result, but
Method B is more flexible (you can cut non-circular holes).

\section{Step 9: Two mounting holes (exercise)}

Using the same approach, cut a ⌀5 through hole on each side. Hints:
\begin{itemize}
  \item draw the circle with the \ui{Circle} tool, add a \ui{Dimension} constraint for the diameter (enter
        \file{5}), then position it with a \ui{Distance Dimension} constraint;
  \item you can also cut one hole first and copy it with \ui{Polar Pattern} or \ui{Linear Pattern};
  \item when you are done the tree should read, in order: Sketch → Pad → Fillet → Hole → Sketch → Pocket →
        Fillet…
\end{itemize}

\section{Step 10: Go back and change parameters}

This is the heart of parametric CAD:

\begin{enumerate}
  \item Double-tap \ui{Pad} in the feature tree → change the length from 12 to 20 → ✓. The whole model updates,
        including the later fillets and holes;
  \item double-tap the first \ui{Sketch} → change 60 to 80 → leave the sketch and the pad grows automatically;
  \item if a later feature fails because of a size change (a hole ending up outside the body, for example),
        the tree shows a red error; double-tap that feature to see why, then change a constraint or a position.
\end{enumerate}

\section{Step 11: Save and export}

\begin{itemize}
  \item \menupath{File → Save}: saves as \file{.FCStd}, FreeCAD's native format, which keeps the whole history;
  \item \menupath{File → Save a copy}: writes a snapshot;
  \item \menupath{File → Export}: choose STEP, IGES, STL, OBJ, 3MF and others for 3D printing or another CAD
        program.
\end{itemize}

\begin{tip}[Advice for 3D printing]
Before exporting STL, check that the solid is valid with \ui{Part → Check Geometry}. STL precision depends on
the “Deviation” parameter: about 0.05 mm for figurine-style models on a tablet, 0.01 mm for mechanical parts.
\end{tip}

\section{Summary}

You have now walked the main trunk of CAD:

\begin{center}
\small
\begin{tabularx}{\linewidth}{@{}lX@{}}
\toprule
\textbf{Concept} & \textbf{What it is in MobileCAD} \\
\midrule
Sketch & A 2D profile plus constraints and dimensions; stable only once fully constrained \\
Pad / Pocket & Add or remove material along the normal of a closed profile \\
Fillet / Chamfer & Dress an edge; lives in the feature history \\
Hole & A standardised hole, with threads and counterbores \\
Feature tree & The modelling history; double-tap any step to change it \\
Tip & Marks how far the body has been computed; roll it back to replay the history \\
\bottomrule
\end{tabularx}
\end{center}

The next chapter covers sketches in depth: every drawing tool, every constraint type, conflict tracking and
B-splines.
